\begin{proposition}[Rank-three period quotients]\hypertarget{tag-01QM}{}\label{fam:src:C.3.1}\label{fam:st:C.3.1}
\StatementRecord{fam:src:C.3.1}{proved}{01QM}{the definitions and constructions stated in the proof}

At every point with \(D=0\), \(T>0\), the real-linear map
\(\Pi:\mathbb R^4\to\mathbb C^2\) has rank three. Write
\(\mu=a_0\tau+b_0\), \(a_0,b_0\in\mathbb R\), and put
\[
 k=(1,-6a_0,-6b_0,6a_0\operatorname{Re}\mu-\operatorname{Re}\beta).
\]
Its real kernel is \(\mathbb R k\). The subgroup \(\Pi\mathbb Z^4\) is
discrete exactly when \(k\in\mathbb Q^4\). In that case it is a
rank-three lattice in a real hyperplane and its quotient in \(\mathbb C^2\)
is a noncompact complex Lie group. In every other case the topological
quotient \(\mathbb C^2/\Pi\mathbb Z^4\) is non-Hausdorff.
\end{proposition}


\begin{proof}
\begin{enumerate}[label=\arabic*.,leftmargin=*]
\item The real matrix, with output order
\((\operatorname{Re}\zeta_1,\operatorname{Im}\zeta_1,
\operatorname{Re}\zeta_2,\operatorname{Im}\zeta_2)\), is
\[
 \begin{pmatrix}6\operatorname{Re}\mu&\operatorname{Re}\tau&1&0\\
 6\operatorname{Im}\mu&T&0&0\\
 \operatorname{Re}\beta&\operatorname{Re}\mu&0&1\\
 \operatorname{Im}\beta&\operatorname{Im}\mu&0&0\end{pmatrix}.
\]
Its last three columns are real-linearly independent because \(T>0\).
Its determinant is \(TD\), so at zero determinant the rank is exactly
three. Substitution gives the displayed kernel, including its signs.
If the image subgroup is discrete, it is a lattice of rank at most three
in its real span. The homomorphism \(\mathbb Z^4\to\Pi\mathbb Z^4\)
therefore has a nonzero kernel, forcing the displayed kernel vector with
first coordinate one to be rational. Conversely a rational line has a
primitive integral generator extendable to a basis of \(\mathbb Z^4\).
The remaining three basis images are a real basis of the image hyperplane
and hence generate a discrete lattice. Its quotient leaves one unbounded
real direction. Finally a non-discrete finitely generated subgroup of
this real hyperplane cannot be closed: a closed subgroup is a product of
a real vector group and a discrete group, and a nonzero real vector group
is uncountable. A quotient topological group is Hausdorff exactly when its
identity subgroup is closed.
\end{enumerate}
\end{proof}